\medskip

\begin{minipage}{0.5\linewidth}
On considère un cube ABCDEFGH et l'espace est rapporté au repère orthonormal
$\left(\text{A}~;~\vect{\text{AB}},~\vect{\text{AD}},~\vect{\text{AE}}\right)$.

Pour tout réel $m$ appartenant à l'intervalle [0~;~1], on considère les points K et L de coordonnées:

\begin{center}K\,$(m~;~0~;~0)$\quad et \quad L\,$(1 - m~;~1~;~1)$.\end{center}
\end{minipage}\hfill
\begin{minipage}{0.45\linewidth}
\psset{unit=0.85cm}
\begin{pspicture}(6.6,6.6)
\pspolygon[fillstyle=solid,fillcolor=lightgray!30,linewidth=0pt](5,5.9)(0.4,5.3)(1.6,0.82)(6.6,1.4)
\pspolygon(0.4,1)(4.2,0.5)(4.2,4.8)(0.4,5.3)
\psline(4.2,0.5)(6.6,1.4)(6.6,5.7)(4.2,4.8)
\psline(6.6,5.7)(2.8,6.2)(0.4,5.3)
\psline(5,5.9)(0.4,5.3)(1.6,0.82)
\psline[linestyle=dashed](5,5.9)(6.6,1.4)(1.6,0.84)
\psline[linestyle=dashed](0.4,1)(2.8,1.9)(6.6,1.4)
\psline[linestyle=dashed](2.8,1.9)(2.8,6.2)
\uput[dl](0.4,1){A} \uput[d](4.2,0.5){B} \uput[dr](6.6,1.4){C} \uput[ul](2.8,1.9){D}
\uput[ul](0.4,5.3){E} \uput[dr](4.2,4.8){F} \uput[ur](6.6,5.7){G} \uput[u](2.8,6.2){H}
\uput[dl](1.6,0.84){K} \uput[u](5,5.9){L}
\end{pspicture}
\end{minipage}

\medskip

\begin{enumerate}
\item Le point E a pour coordonnées (0~;~0~;~1).

Comme $\vectt{AC}=\vectt{AB}+\vectt{AD}$, le point C a pour coordonnées (1~;~1~;~0).

\item Dans cette question, $m = 0$. Ainsi, le point L\,(1~;~1~;~1) est confondu avec le point G, le point K(0~;~0~;~0) est confondu avec le point A et le plan (LEK) est donc le plan (GEA).
	\begin{enumerate}
		\item 
\begin{list}{\textbullet}{}
\item ABCD est un carré donc ses diagonales sont perpendiculaires, donc $\vectt{DB} \perp \vectt{AC}$.
\item Le vecteur $\vectt{AE}$ est normal au plan (ABD) donc $\vectt{AE}$ est orthogonal à tout vecteur du plan (ABD), donc $\vectt{AE} \perp \vectt{DB}$.
\item Les vecteurs $\vectt{AE}$ et $\vectt{AC}$ ne sont pas colinéaires. 
\end{list}		
		
Le vecteur $\vectt{DB}$	est orthogonal à deux vecteurs non colinéaires du plan (GEA) donc le vecteur $\vectt{DB}$ est normal au plan (GEA).
		
		\item 
On déduit de la question précédente que le plan (GEA) a une équation de la forme $1\times x + (-1)\times y + 0\times z + d =0$, soit $x-y+d=0$.

Le plan (GEA) passe par le point A de coordonnées (0~;~0~;~0) donc $0-0+d=0$ et donc $d=0$.

Le plan (GEA) a donc pour équation $x-y=0$.
	\end{enumerate}
\end{enumerate}

On s'intéresse désormais à la nature de CKEL en fonction du paramètre $m$.

\begin{enumerate}[resume]
\item Dans cette question, $m$ est un réel quelconque de l'intervalle [0~;~1].
	\begin{enumerate}
		\item 
\begin{list}{\textbullet}{}
\item 
K\,$\begin{pmatrix} m\\0\\0 \end{pmatrix}$ et
E\,$\begin{pmatrix} 0\\0\\1 \end{pmatrix}$
donc $\vectt{KE}$ a pour coordonnées
$\begin{pmatrix} 0-m\\0-0\\1-0 \end{pmatrix} = \begin{pmatrix} -m\\0\\1 \end{pmatrix}$.
\item 
C\,$\begin{pmatrix} 1\\1\\0 \end{pmatrix}$ et
L\,$\begin{pmatrix} 1-m\\1\\1 \end{pmatrix}$ 
donc $\vectt{CL}$ a pour coordonnées
$\begin{pmatrix} 1-m-1\\1-1\\1-0 \end{pmatrix} = \begin{pmatrix} -m\\0\\1 \end{pmatrix}$.
\end{list}		

$\vectt{KE} = \vectt{CL}$ donc CKEL est un parallélogramme.
		
		\item
K\,$\begin{pmatrix} m\\0\\0 \end{pmatrix}$ et
C\,$\begin{pmatrix} 1\\1\\0 \end{pmatrix}$
donc $\vectt{KC}$ a pour coordonnées
$\begin{pmatrix} 1-m\\1-0\\0-0 \end{pmatrix} = \begin{pmatrix} 1-m\\1\\0 \end{pmatrix}$.

$\vectt{KC} \cdot \vectt{KE} 
= x_{\vectt{KC}}\times x_{\vectt{KE}} + y_{\vectt{KC}}\times y_{\vectt{KE}}  + z_{\vectt{KC}}\times z_{\vectt{KE}} 
= (1-m)\times (-m) + 1 \times 0 + 0 \times (-1)
= m(m - 1)$	
		
		\item
\begin{tabular}[t]{@{} l @{ si, et seulement si, } l}
CKEL est un rectangle & CKEL possède un angle droit\\
&$\vectt{KC} \perp \vectt{KE}$\\
& $\vectt{KC} \cdot \vectt{KE}=0$\\
& $m(m-1)=0$\\
&$m = 0$ ou $m = 1$
\end{tabular}	

\end{enumerate}
\item Dans cette question, $m = \dfrac12$. 

Ainsi, L a pour coordonnées  $\left(\dfrac12~;~1~;~1\right)$ et K a pour coordonnées $\left(\dfrac12~;~0~;~0\right)$.
	\begin{enumerate}
		\item 
\begin{list}{\textbullet}{}
\item
K\,$\begin{pmatrix} \frac{1}{2}\\0\\0 \end{pmatrix}$ et
L\,$\begin{pmatrix} \frac{1}{2}\\1\\1 \end{pmatrix}$
donc $\vectt{KL}$ a pour coordonnées
$\begin{pmatrix} \frac{1}{2}-\frac{1}{2}\\1-0\\1-0 \end{pmatrix} = \begin{pmatrix} 0\\1\\1 \end{pmatrix}$.		
\item
E\,$\begin{pmatrix} 0\\0\\1 \end{pmatrix}$ et
C\,$\begin{pmatrix} 1\\1\\0  \end{pmatrix}$
donc $\vectt{EC}$ a pour coordonnées
$\begin{pmatrix} 1-0\\1-0\\0-1 \end{pmatrix} = \begin{pmatrix} 1\\1\\-1 \end{pmatrix}$.		
\item
$\vectt{KL} \cdot \vectt{EC} = 0\times 1 + 1\times 1 + 1 \times (-1)=0$ donc $\vectt{KL} \perp \vectt{EC}$
\end{list}
		
Le parallélogramme CKEL a ses diagonales perpendiculaires donc c'est un losange.		
		
		\item 
\begin{list}{\textbullet}{}
\item $\vectt{KC} \cdot \vectt{KE}  = m(m-1)$; or $m=\dfrac{1}{2}$ donc $\vectt{KC} \cdot \vectt{KE} =\dfrac{1}{2} \left ( \dfrac{1}{2}-1 \right ) = - \dfrac{1}{4}$
\item $\vectt{KC} \cdot \vectt{KE}  = \text{KC} \times \text{KE} \times \cos \left ( \widehat{\text{CKE}} \right )$

\item $\vectt{KC}$ a pour coordonnées $\begin{pmatrix} 1-m\\1\\0 \end{pmatrix} = \begin{pmatrix}  \frac{1}{2}\\1\\0 \end{pmatrix}$

Donc $\text{KC}^2 = \left ( \dfrac{1}{2} \right )^2+ 1^2+0^2=\dfrac{5}{4}$ donc $\text{KC}=\dfrac{\sqrt{5}}{2}$

\item $\vectt{KE}$ a pour coordonnées $\begin{pmatrix} -m\\0\\1 \end{pmatrix} = \begin{pmatrix}  -\frac{1}{2}\\0\\1 \end{pmatrix}$

Donc $\text{KE}^2 = \left (-\dfrac{1}{2} \right )^2+ 0^2+1^2=\dfrac{5}{4}$ donc $\text{KE}=\dfrac{\sqrt{5}}{2}$
\end{list}		

On a donc:
$-\dfrac{1}{4} = \dfrac{\sqrt{5}}{2}\times \dfrac{\sqrt{5}}{2} \times \cos \left ( \widehat{\text{CKE}} \right )$

Donc $-\dfrac{1}{4} = \dfrac{5}{4}\times \cos \left ( \widehat{\text{CKE}} \right )$ et donc  $\cos \left ( \widehat{\text{CKE}} \right ) = -\dfrac{1}{5}$

On en déduit avec la calculatrice que $\widehat{\text{CKE}} \approx 101,5\degres$, soit $\widehat{\text{CKE}} \approx 102\degres$.
	\end{enumerate}
\end{enumerate}
